\begin{frame}{\insertsubsection}

  CVE-2017-5715: тренировка предсказателя переходов

  \begin{figure}[h]
    \includegraphics[height = .7\textheight]{spectre_logo}
    % \caption{Spectre}
  \end{figure}

  \note{


  }
\end{frame}

% \subsubsection{Предсказатель переходов и его тренировка}
% \begin{frame}{\insertsubsubsection}

%   \begin{figure}[h]
%     \begin{tikzpicture}[
%       align = center,
%       ->,
%       > = Stealth,
%       thick,
%       ampersand replacement = \&,
%       block/.style = {
%         rectangle,
%         draw,
%         fill = ForestGreen,
%         text = White,
%         text centered,
%       },
%       ]

%       \node[
%       text = Black,
%       font = \ttfamily,
%       ] (code) {
%         Animal* a = \color{ForestGreen}\only<1-4>{bird}\only<5->{fish};
%       };

%       \node[
%       below = .5 of code,
%       text = Black,
%       font = \ttfamily,
%       ] (move) {
%         a->move()
%       };

%       \node[
%       below left = 3 of move,
%       text = Black,
%       font = \ttfamily,
%       minimum width = 3cm,
%       ] (cache) {
%         arr1[untrusted]
%       };

%       \node[
%       circle,
%       below = of move,
%       text = White,
%       fill = ForestGreen,
%       draw,
%       label = below:Предсказание,
%       ] (zero) {
%         \only<1-3>{swim()}%
%         \only<4->{fly()}%
%       };

%       \node[
%       below right = 3 of move,
%       text = Black,
%       font = \ttfamily,
%       minimum width = 3cm,
%       ] (zero) {
%         0
%       };

%       \draw<1-2,4-5,7> (move) -> node[above, sloped] {fly()} (cache);
%       \draw<1,3,4-6> (move) -> node[above, sloped] {swim()} (zero);

%       \draw<2>[ForestGreen] (move) -> node[above, sloped] {swim()} node[below, sloped] {Спекулятивно} (zero);
%       \draw<3>[ForestGreen] (move) -> node[above, sloped] {fly()} node[below, sloped] {Реально} (cache);

%       \draw<6>[ForestGreen] (move) -> node[above, sloped] {fly()} node[below, sloped] {Спекулятивно} (cache);
%       \draw<7>[ForestGreen] (move) -> node[above, sloped] {swim()} node[below, sloped] {Реально} (zero);

%     \end{tikzpicture}
%     % \caption{}

%   \end{figure}

%   \note<1>{

%     Ранее уже рассказывалось про предсказатель переходов, а также про буфер.

%     Для атаки требуется досконально знать, \textbf{как работает предсказатель
%       переходов}.

%   }

%   \note<6>{

%     \textbf{Внимание}, спекулятивно выполняется натренированная нами ветка! В
%     итоге будет исполняться код, который совершит атаку на кэш.

%   }
% \end{frame}

% \begin{frame}{\insertsubsubsection}

%   \begin{figure}[h]
%     \begin{tikzpicture}[
%       align = center,
%       ->,
%       > = Stealth,
%       thick,
%       ampersand replacement = \&,
%       bp/.style = {
%         rectangle,
%         rectangle split,
%         rectangle split parts = 2,
%         rectangle split horizontal,
%         rectangle split part fill = {
%           ForestGreen, NavyBlue
%         },
%         draw,
%         text = White,
%       },
%       ]

%       \node[
%       bp,
%       ] (bp) {
%         \color{Black}0xEBE45A82
%         \nodepart{two}
%         T,T,N,N,T,T,N,N
%         \nodepart{three}
%       };

%       \node[
%       bp,
%       above left = of bp.center,
%       label = above:Процесс A,
%       ] (proc_a) {
%         0x0000 \color{Black}EBE45A82
%         \nodepart{two}
%         Переход A
%         \nodepart{three}
%       };

%       \node[
%       bp,
%       above right = of bp.center,
%       label = above:Процесс B (ядро/гипервизор),
%       ] (proc_b) {
%         0xFFFF \color{Black}EBE45A82
%         \nodepart{two}
%         Переход B
%         \nodepart{three}
%       };

%       \draw (proc_a.one south) |- (bp.one west);
%       \draw (proc_b.one west) -> (bp.one north);


%     \end{tikzpicture}
%     \caption{В BTB используются виртуальные адреса, а также возникают коллизии}

%   \end{figure}


%   \note{

%     Из-за того, что возникают коллизии в таблице BTB, мы можем
%     \textbf{натренировать} его таким образом, чтобы спекулятивно исполнялся
%     нужный нам переход.

%     Существует \textbf{не один способ} тренировки предсказателя переходов.

%     Также существует возможность создания ROP цепочки из \textbf{гаджетов
%       программы-жертвы} и натренировать на него, но для этого требуется знать
%     адрес. Чтобы \textbf{узнать адрес перехода} можно применить \textbf{атаку по
%       сторонним каналам} на предсказатель переходов.

%   }
% \end{frame}

% \subsubsection{И снова спекулятивное выполнение}
% \begin{frame}[fragile]{\insertsubsubsection}

%   \begin{minted}{c}
%   if (untrusted_offset_from_user < array1_size)
%     y = array2[((array1[untrusted_offset_from_user] & 1) * 0x100) + 0x200];
%   \end{minted}

%   \note{

%     Ничего нового, используется всё тот же код, чаще всего ROP цепочка, которая
%     приводит к атакам на кэш.

%   }
% \end{frame}

\subsubsection{Предотвращение}
\begin{frame}{\insertsubsubsection}

  Частично такое же, как и в случае с Variant 1

  \begin{itemize}
  \item отключение предсказателя переходов
  \item очистка буфера предсказателя переходов при переключении контекста
  \item вставка «барьеров» (Indirect Branch Restrict Speculation, Indirect
    Branch Predictor Barrier и др.)
  \item retpoline --- «оборачивание» косвенных переходов
  \item oo7 --- умное «оборачивание» косвенных переходов
  \end{itemize}

  \note{

    Всё медленно!

    retpoline --- замена всех косвенных переходов, дополнение инструкций
    возврата, паузы перед непрямыми вызовами функций

  }

\end{frame}